% Load with % Load with % Load with % Load with \input{macros.tex} before \begin{document}.
% Edit or remove any definition below to suit your paper.
% Some notation follows Truong-Son Van's GiantCluster macros
% (originally credited to Gautam Iyer).

% Standard packages; no custom style files are needed.
\usepackage{amsmath,amssymb,amsthm}
\usepackage{mathtools}
\usepackage{xcolor}
\usepackage[mathlines]{lineno}
\usepackage[hidelinks]{hyperref}

% Line numbers: use \linenumbers in the document to enable them,
% \nolinenumbers to stop, or \modulolinenumbers[5] to label every fifth line.

% Named sidenotes: \sidenote[Son]{Check this argument.}
% Use \shownotesfalse in the preamble to hide all sidenotes.
\newif\ifshownotes
\shownotestrue
\setlength{\marginparwidth}{1in}
\setlength{\marginparsep}{10pt}
\newcommand{\sidenote}[2][Note]{%
  \ifshownotes
    \marginpar{\raggedright\footnotesize\color{blue!60!black}%
      \textbf{#1:} #2}%
  \fi
}
% Optional shortcut for a collaborator (copy and rename as needed):
% \newcommand{\sonnote}[1]{\sidenote[Son]{#1}}

% Theorem environments and equation numbers, numbered by section.
\numberwithin{equation}{section}
\theoremstyle{plain}
\newtheorem{theorem}{Theorem}[section]
\newtheorem{lemma}[theorem]{Lemma}
\newtheorem{proposition}[theorem]{Proposition}
\newtheorem{corollary}[theorem]{Corollary}
\newtheorem{conjecture}[theorem]{Conjecture}
\theoremstyle{definition}
\newtheorem{definition}[theorem]{Definition}
\newtheorem{example}[theorem]{Example}
\theoremstyle{remark}
\newtheorem{remark}[theorem]{Remark}

% Number systems.
\newcommand{\R}{\mathbb{R}}
\newcommand{\C}{\mathbb{C}}
\newcommand{\N}{\mathbb{N}}
\newcommand{\Z}{\mathbb{Z}}
\newcommand{\Q}{\mathbb{Q}}
\newcommand{\T}{\mathbb{T}}

% Calligraphic letters.
\newcommand{\cA}{\mathcal{A}}
\newcommand{\cB}{\mathcal{B}}
\newcommand{\cC}{\mathcal{C}}
\newcommand{\cD}{\mathcal{D}}
\newcommand{\cE}{\mathcal{E}}
\newcommand{\cF}{\mathcal{F}}
\newcommand{\cH}{\mathcal{H}}
\newcommand{\cL}{\mathcal{L}}
\newcommand{\cM}{\mathcal{M}}
\newcommand{\cP}{\mathcal{P}}
\newcommand{\cQ}{\mathcal{Q}}
\newcommand{\cV}{\mathcal{V}}
\newcommand{\cX}{\mathcal{X}}
\newcommand{\cY}{\mathcal{Y}}

% Delimiters: \norm{u}, \norm*{...} for automatic sizing,
% or \norm[\big]{u} for an explicit size. The same applies to the others.
\DeclarePairedDelimiter{\abs}{\lvert}{\rvert}
\DeclarePairedDelimiter{\norm}{\lVert}{\rVert}
\DeclarePairedDelimiter{\paren}{(}{)}
\DeclarePairedDelimiter{\brak}{[}{]}
\DeclarePairedDelimiter{\set}{\lbrace}{\rbrace}
\DeclarePairedDelimiter{\av}{\langle}{\rangle}
\DeclarePairedDelimiter{\floor}{\lfloor}{\rfloor}
\DeclarePairedDelimiter{\ceil}{\lceil}{\rceil}
\DeclarePairedDelimiterX{\ip}[2]{\langle}{\rangle}{#1,#2}
\newcommand{\st}{\mid}
\newcommand{\given}{\mid}

% Calculus. Optional argument = order of differentiation.
\newcommand{\dd}{\mathop{}\!\mathrm{d}}
\newcommand{\D}{\partial}
\newcommand{\odv}[3][]{\frac{\mathrm{d}^{#1}#2}{\mathrm{d}#3^{#1}}}
\newcommand{\pdv}[3][]{\frac{\partial^{#1}#2}{\partial #3^{#1}}}
\newcommand{\grad}{\nabla}
\newcommand{\dv}{\nabla\cdot}
\newcommand{\curl}{\nabla\times}
\newcommand{\lap}{\Delta}

% Probability. Use \Prob so the standard \P symbol remains available.
\newcommand{\E}{\mathbb{E}}
\newcommand{\Prob}{\mathbb{P}}
\newcommand{\one}{\mathbf{1}}
\DeclareMathOperator{\var}{Var}
\DeclareMathOperator{\Cov}{Cov}

% Common operators.
\DeclareMathOperator{\supp}{supp}
\DeclareMathOperator{\dist}{dist}
\DeclareMathOperator{\diam}{diam}
\DeclareMathOperator{\trace}{tr}
\DeclareMathOperator{\rank}{rank}
\DeclareMathOperator{\Span}{span}
\DeclareMathOperator{\Lip}{Lip}
\DeclareMathOperator*{\esssup}{ess\,sup}
\DeclareMathOperator*{\argmin}{arg\,min}
\DeclareMathOperator*{\argmax}{arg\,max}

% Miscellaneous notation.
\newcommand{\eps}{\varepsilon}
\newcommand{\id}{\mathrm{id}}
\newcommand{\defeq}{\coloneqq}
\newcommand{\weakto}{\rightharpoonup}
\newcommand{\weakstar}{\overset{\ast}{\rightharpoonup}}

% Add notation specific to your document below this line.
 before \begin{document}.
% Edit or remove any definition below to suit your paper.
% Some notation follows Truong-Son Van's GiantCluster macros
% (originally credited to Gautam Iyer).

% Standard packages; no custom style files are needed.
\usepackage{amsmath,amssymb,amsthm}
\usepackage{mathtools}
\usepackage{xcolor}
\usepackage[mathlines]{lineno}
\usepackage[hidelinks]{hyperref}

% Line numbers: use \linenumbers in the document to enable them,
% \nolinenumbers to stop, or \modulolinenumbers[5] to label every fifth line.

% Named sidenotes: \sidenote[Son]{Check this argument.}
% Use \shownotesfalse in the preamble to hide all sidenotes.
\newif\ifshownotes
\shownotestrue
\setlength{\marginparwidth}{1in}
\setlength{\marginparsep}{10pt}
\newcommand{\sidenote}[2][Note]{%
  \ifshownotes
    \marginpar{\raggedright\footnotesize\color{blue!60!black}%
      \textbf{#1:} #2}%
  \fi
}
% Optional shortcut for a collaborator (copy and rename as needed):
% \newcommand{\sonnote}[1]{\sidenote[Son]{#1}}

% Theorem environments and equation numbers, numbered by section.
\numberwithin{equation}{section}
\theoremstyle{plain}
\newtheorem{theorem}{Theorem}[section]
\newtheorem{lemma}[theorem]{Lemma}
\newtheorem{proposition}[theorem]{Proposition}
\newtheorem{corollary}[theorem]{Corollary}
\newtheorem{conjecture}[theorem]{Conjecture}
\theoremstyle{definition}
\newtheorem{definition}[theorem]{Definition}
\newtheorem{example}[theorem]{Example}
\theoremstyle{remark}
\newtheorem{remark}[theorem]{Remark}

% Number systems.
\newcommand{\R}{\mathbb{R}}
\newcommand{\C}{\mathbb{C}}
\newcommand{\N}{\mathbb{N}}
\newcommand{\Z}{\mathbb{Z}}
\newcommand{\Q}{\mathbb{Q}}
\newcommand{\T}{\mathbb{T}}

% Calligraphic letters.
\newcommand{\cA}{\mathcal{A}}
\newcommand{\cB}{\mathcal{B}}
\newcommand{\cC}{\mathcal{C}}
\newcommand{\cD}{\mathcal{D}}
\newcommand{\cE}{\mathcal{E}}
\newcommand{\cF}{\mathcal{F}}
\newcommand{\cH}{\mathcal{H}}
\newcommand{\cL}{\mathcal{L}}
\newcommand{\cM}{\mathcal{M}}
\newcommand{\cP}{\mathcal{P}}
\newcommand{\cQ}{\mathcal{Q}}
\newcommand{\cV}{\mathcal{V}}
\newcommand{\cX}{\mathcal{X}}
\newcommand{\cY}{\mathcal{Y}}

% Delimiters: \norm{u}, \norm*{...} for automatic sizing,
% or \norm[\big]{u} for an explicit size. The same applies to the others.
\DeclarePairedDelimiter{\abs}{\lvert}{\rvert}
\DeclarePairedDelimiter{\norm}{\lVert}{\rVert}
\DeclarePairedDelimiter{\paren}{(}{)}
\DeclarePairedDelimiter{\brak}{[}{]}
\DeclarePairedDelimiter{\set}{\lbrace}{\rbrace}
\DeclarePairedDelimiter{\av}{\langle}{\rangle}
\DeclarePairedDelimiter{\floor}{\lfloor}{\rfloor}
\DeclarePairedDelimiter{\ceil}{\lceil}{\rceil}
\DeclarePairedDelimiterX{\ip}[2]{\langle}{\rangle}{#1,#2}
\newcommand{\st}{\mid}
\newcommand{\given}{\mid}

% Calculus. Optional argument = order of differentiation.
\newcommand{\dd}{\mathop{}\!\mathrm{d}}
\newcommand{\D}{\partial}
\newcommand{\odv}[3][]{\frac{\mathrm{d}^{#1}#2}{\mathrm{d}#3^{#1}}}
\newcommand{\pdv}[3][]{\frac{\partial^{#1}#2}{\partial #3^{#1}}}
\newcommand{\grad}{\nabla}
\newcommand{\dv}{\nabla\cdot}
\newcommand{\curl}{\nabla\times}
\newcommand{\lap}{\Delta}

% Probability. Use \Prob so the standard \P symbol remains available.
\newcommand{\E}{\mathbb{E}}
\newcommand{\Prob}{\mathbb{P}}
\newcommand{\one}{\mathbf{1}}
\DeclareMathOperator{\var}{Var}
\DeclareMathOperator{\Cov}{Cov}

% Common operators.
\DeclareMathOperator{\supp}{supp}
\DeclareMathOperator{\dist}{dist}
\DeclareMathOperator{\diam}{diam}
\DeclareMathOperator{\trace}{tr}
\DeclareMathOperator{\rank}{rank}
\DeclareMathOperator{\Span}{span}
\DeclareMathOperator{\Lip}{Lip}
\DeclareMathOperator*{\esssup}{ess\,sup}
\DeclareMathOperator*{\argmin}{arg\,min}
\DeclareMathOperator*{\argmax}{arg\,max}

% Miscellaneous notation.
\newcommand{\eps}{\varepsilon}
\newcommand{\id}{\mathrm{id}}
\newcommand{\defeq}{\coloneqq}
\newcommand{\weakto}{\rightharpoonup}
\newcommand{\weakstar}{\overset{\ast}{\rightharpoonup}}

% Add notation specific to your document below this line.
 before \begin{document}.
% Edit or remove any definition below to suit your paper.
% Some notation follows Truong-Son Van's GiantCluster macros
% (originally credited to Gautam Iyer).

% Standard packages; no custom style files are needed.
\usepackage{amsmath,amssymb,amsthm}
\usepackage{mathtools}
\usepackage{xcolor}
\usepackage[mathlines]{lineno}
\usepackage[hidelinks]{hyperref}

% Line numbers: use \linenumbers in the document to enable them,
% \nolinenumbers to stop, or \modulolinenumbers[5] to label every fifth line.

% Named sidenotes: \sidenote[Son]{Check this argument.}
% Use \shownotesfalse in the preamble to hide all sidenotes.
\newif\ifshownotes
\shownotestrue
\setlength{\marginparwidth}{1in}
\setlength{\marginparsep}{10pt}
\newcommand{\sidenote}[2][Note]{%
  \ifshownotes
    \marginpar{\raggedright\footnotesize\color{blue!60!black}%
      \textbf{#1:} #2}%
  \fi
}
% Optional shortcut for a collaborator (copy and rename as needed):
% \newcommand{\sonnote}[1]{\sidenote[Son]{#1}}

% Theorem environments and equation numbers, numbered by section.
\numberwithin{equation}{section}
\theoremstyle{plain}
\newtheorem{theorem}{Theorem}[section]
\newtheorem{lemma}[theorem]{Lemma}
\newtheorem{proposition}[theorem]{Proposition}
\newtheorem{corollary}[theorem]{Corollary}
\newtheorem{conjecture}[theorem]{Conjecture}
\theoremstyle{definition}
\newtheorem{definition}[theorem]{Definition}
\newtheorem{example}[theorem]{Example}
\theoremstyle{remark}
\newtheorem{remark}[theorem]{Remark}

% Number systems.
\newcommand{\R}{\mathbb{R}}
\newcommand{\C}{\mathbb{C}}
\newcommand{\N}{\mathbb{N}}
\newcommand{\Z}{\mathbb{Z}}
\newcommand{\Q}{\mathbb{Q}}
\newcommand{\T}{\mathbb{T}}

% Calligraphic letters.
\newcommand{\cA}{\mathcal{A}}
\newcommand{\cB}{\mathcal{B}}
\newcommand{\cC}{\mathcal{C}}
\newcommand{\cD}{\mathcal{D}}
\newcommand{\cE}{\mathcal{E}}
\newcommand{\cF}{\mathcal{F}}
\newcommand{\cH}{\mathcal{H}}
\newcommand{\cL}{\mathcal{L}}
\newcommand{\cM}{\mathcal{M}}
\newcommand{\cP}{\mathcal{P}}
\newcommand{\cQ}{\mathcal{Q}}
\newcommand{\cV}{\mathcal{V}}
\newcommand{\cX}{\mathcal{X}}
\newcommand{\cY}{\mathcal{Y}}

% Delimiters: \norm{u}, \norm*{...} for automatic sizing,
% or \norm[\big]{u} for an explicit size. The same applies to the others.
\DeclarePairedDelimiter{\abs}{\lvert}{\rvert}
\DeclarePairedDelimiter{\norm}{\lVert}{\rVert}
\DeclarePairedDelimiter{\paren}{(}{)}
\DeclarePairedDelimiter{\brak}{[}{]}
\DeclarePairedDelimiter{\set}{\lbrace}{\rbrace}
\DeclarePairedDelimiter{\av}{\langle}{\rangle}
\DeclarePairedDelimiter{\floor}{\lfloor}{\rfloor}
\DeclarePairedDelimiter{\ceil}{\lceil}{\rceil}
\DeclarePairedDelimiterX{\ip}[2]{\langle}{\rangle}{#1,#2}
\newcommand{\st}{\mid}
\newcommand{\given}{\mid}

% Calculus. Optional argument = order of differentiation.
\newcommand{\dd}{\mathop{}\!\mathrm{d}}
\newcommand{\D}{\partial}
\newcommand{\odv}[3][]{\frac{\mathrm{d}^{#1}#2}{\mathrm{d}#3^{#1}}}
\newcommand{\pdv}[3][]{\frac{\partial^{#1}#2}{\partial #3^{#1}}}
\newcommand{\grad}{\nabla}
\newcommand{\dv}{\nabla\cdot}
\newcommand{\curl}{\nabla\times}
\newcommand{\lap}{\Delta}

% Probability. Use \Prob so the standard \P symbol remains available.
\newcommand{\E}{\mathbb{E}}
\newcommand{\Prob}{\mathbb{P}}
\newcommand{\one}{\mathbf{1}}
\DeclareMathOperator{\var}{Var}
\DeclareMathOperator{\Cov}{Cov}

% Common operators.
\DeclareMathOperator{\supp}{supp}
\DeclareMathOperator{\dist}{dist}
\DeclareMathOperator{\diam}{diam}
\DeclareMathOperator{\trace}{tr}
\DeclareMathOperator{\rank}{rank}
\DeclareMathOperator{\Span}{span}
\DeclareMathOperator{\Lip}{Lip}
\DeclareMathOperator*{\esssup}{ess\,sup}
\DeclareMathOperator*{\argmin}{arg\,min}
\DeclareMathOperator*{\argmax}{arg\,max}

% Miscellaneous notation.
\newcommand{\eps}{\varepsilon}
\newcommand{\id}{\mathrm{id}}
\newcommand{\defeq}{\coloneqq}
\newcommand{\weakto}{\rightharpoonup}
\newcommand{\weakstar}{\overset{\ast}{\rightharpoonup}}

% Add notation specific to your document below this line.
 before \begin{document}.
% Edit or remove any definition below to suit your paper.
% Some notation follows Truong-Son Van's GiantCluster macros
% (originally credited to Gautam Iyer).

% Standard packages; no custom style files are needed.
\usepackage{amsmath,amssymb,amsthm}
\usepackage{mathtools}
\usepackage{xcolor}
\usepackage[mathlines]{lineno}
\usepackage[hidelinks]{hyperref}

% Line numbers: use \linenumbers in the document to enable them,
% \nolinenumbers to stop, or \modulolinenumbers[5] to label every fifth line.

% Named sidenotes: \sidenote[Son]{Check this argument.}
% Use \shownotesfalse in the preamble to hide all sidenotes.
\newif\ifshownotes
\shownotestrue
\setlength{\marginparwidth}{1in}
\setlength{\marginparsep}{10pt}
\newcommand{\sidenote}[2][Note]{%
  \ifshownotes
    \marginpar{\raggedright\footnotesize\color{blue!60!black}%
      \textbf{#1:} #2}%
  \fi
}
% Optional shortcut for a collaborator (copy and rename as needed):
% \newcommand{\sonnote}[1]{\sidenote[Son]{#1}}

% Theorem environments and equation numbers, numbered by section.
\numberwithin{equation}{section}
\theoremstyle{plain}
\newtheorem{theorem}{Theorem}[section]
\newtheorem{lemma}[theorem]{Lemma}
\newtheorem{proposition}[theorem]{Proposition}
\newtheorem{corollary}[theorem]{Corollary}
\newtheorem{conjecture}[theorem]{Conjecture}
\theoremstyle{definition}
\newtheorem{definition}[theorem]{Definition}
\newtheorem{example}[theorem]{Example}
\theoremstyle{remark}
\newtheorem{remark}[theorem]{Remark}

% Number systems.
\newcommand{\R}{\mathbb{R}}
\newcommand{\C}{\mathbb{C}}
\newcommand{\N}{\mathbb{N}}
\newcommand{\Z}{\mathbb{Z}}
\newcommand{\Q}{\mathbb{Q}}
\newcommand{\T}{\mathbb{T}}

% Calligraphic letters.
\newcommand{\cA}{\mathcal{A}}
\newcommand{\cB}{\mathcal{B}}
\newcommand{\cC}{\mathcal{C}}
\newcommand{\cD}{\mathcal{D}}
\newcommand{\cE}{\mathcal{E}}
\newcommand{\cF}{\mathcal{F}}
\newcommand{\cH}{\mathcal{H}}
\newcommand{\cL}{\mathcal{L}}
\newcommand{\cM}{\mathcal{M}}
\newcommand{\cP}{\mathcal{P}}
\newcommand{\cQ}{\mathcal{Q}}
\newcommand{\cV}{\mathcal{V}}
\newcommand{\cX}{\mathcal{X}}
\newcommand{\cY}{\mathcal{Y}}

% Delimiters: \norm{u}, \norm*{...} for automatic sizing,
% or \norm[\big]{u} for an explicit size. The same applies to the others.
\DeclarePairedDelimiter{\abs}{\lvert}{\rvert}
\DeclarePairedDelimiter{\norm}{\lVert}{\rVert}
\DeclarePairedDelimiter{\paren}{(}{)}
\DeclarePairedDelimiter{\brak}{[}{]}
\DeclarePairedDelimiter{\set}{\lbrace}{\rbrace}
\DeclarePairedDelimiter{\av}{\langle}{\rangle}
\DeclarePairedDelimiter{\floor}{\lfloor}{\rfloor}
\DeclarePairedDelimiter{\ceil}{\lceil}{\rceil}
\DeclarePairedDelimiterX{\ip}[2]{\langle}{\rangle}{#1,#2}
\newcommand{\st}{\mid}
\newcommand{\given}{\mid}

% Calculus. Optional argument = order of differentiation.
\newcommand{\dd}{\mathop{}\!\mathrm{d}}
\newcommand{\D}{\partial}
\newcommand{\odv}[3][]{\frac{\mathrm{d}^{#1}#2}{\mathrm{d}#3^{#1}}}
\newcommand{\pdv}[3][]{\frac{\partial^{#1}#2}{\partial #3^{#1}}}
\newcommand{\grad}{\nabla}
\newcommand{\dv}{\nabla\cdot}
\newcommand{\curl}{\nabla\times}
\newcommand{\lap}{\Delta}

% Probability. Use \Prob so the standard \P symbol remains available.
\newcommand{\E}{\mathbb{E}}
\newcommand{\Prob}{\mathbb{P}}
\newcommand{\one}{\mathbf{1}}
\DeclareMathOperator{\var}{Var}
\DeclareMathOperator{\Cov}{Cov}

% Common operators.
\DeclareMathOperator{\supp}{supp}
\DeclareMathOperator{\dist}{dist}
\DeclareMathOperator{\diam}{diam}
\DeclareMathOperator{\trace}{tr}
\DeclareMathOperator{\rank}{rank}
\DeclareMathOperator{\Span}{span}
\DeclareMathOperator{\Lip}{Lip}
\DeclareMathOperator*{\esssup}{ess\,sup}
\DeclareMathOperator*{\argmin}{arg\,min}
\DeclareMathOperator*{\argmax}{arg\,max}

% Miscellaneous notation.
\newcommand{\eps}{\varepsilon}
\newcommand{\id}{\mathrm{id}}
\newcommand{\defeq}{\coloneqq}
\newcommand{\weakto}{\rightharpoonup}
\newcommand{\weakstar}{\overset{\ast}{\rightharpoonup}}

% Add notation specific to your document below this line.
